\subsection{A resolved coefficient-tail quadratic repair on a sphere parameter}

We use $S_R=\Gamma_R(T)$ for $R=\mathbb Z$ or
$\mathbb Z_{(p)}$, and the shift convention of AT01--AT05.
In particular
\[
 R\operatorname{Hom}_{\mathrm{gr}S_R}
       (I_N\langle s\rangle,S_R)=S_R^s[0]\qquad(s\geq0),
 \tag{AT.71}
\]
with zero positive Ext. A positive source shift does
not make its degree-zero Hom vanish.
For $R=\mathbb Z$ this is AT04.
For $R=\mathbb Z_{(p)}$ the same assertion follows from
$0\to S_R\to S_{\mathbb Q}\oplus S_{\mathbb Z_p}
\to S_{\mathbb Q_p}\to0$:
the completed-base tail comparison is Theorem 8.3 of
Nagpal--Snowden, as applied in AT02;
the rationalized tail is free, and the degree-zero
difference map is $\mathbb Q\oplus\mathbb Z_p\to\mathbb Q_p$,
which is surjective. The Hom comparison in degree zero
is restriction from the coefficient copy, as before.

For an admissible target $A$, set
\[
 C_A^{w,j}
   =\operatorname{Hom}_{s\operatorname{Mod}_R}
           (K(2),A^{w,j})
   =Z_2\bigl(N(A^{w,j})\bigr).
 \tag{AT.72}
\]
This is a cochain complex in $j$, with its coefficient
$S_R$-action. For weights above the horizontal and
vertical bounds through degree $d$, its cohomology
through degree $d$ is $S_R^w$ in degree two and zero
in the other degrees.

Here is the complete comparison argument for the range
used below. Uniform vertical comparison makes the
cohomology of $A_0$ equal to $S_R^w$ in degree zero
and zero in positive degrees through $d$.
The exact normalization functors make $N_1A,N_2A$
acyclic through $d$, since positive normalization of
the constant coefficient module is zero.
Write $Z_1=\ker(d_0:N_1A\to A_0)$ and
$B_1=\operatorname{im}(d_0:N_2A\to N_1A)$.
The actual short exact cochain sequences
\[
 0\to Z_1\to N_1A\to\operatorname{im}(N_1A\to A_0)\to0,
 \qquad
 0\to C_A\to N_2A\to B_1\to0
 \tag{AT.73}
\]
give the assertion: horizontal exactness identifies
the first image with $A_0$ and $B_1$ with $Z_1$
in every form degree at most $d$.
To compute $H^j(C_A)$ with $j\leq d$, its long exact
sequence uses only $H^{j-1}(B_1)$ and $H^j(N_2A)$;
the former uses terms through degree $j$.
Thus no assertion of row exactness beyond the stated
finite range is hidden here.
The resulting period identification is natural and
$S_R$-linear. No strict $S_R$-linear choice of coefficient
edge primitives has been assumed.

\begin{theorem}[AT16: Proved; an actual resolved-tail quadratic repair]
\label{AT16}
Let $X\to A,B$ be actual strict simplicial PD-DGA maps
over $R$ with torsionfree terms, with $A,B$ admissible.
Suppose $H^0(X_a)=S_R$ in every weight and simplex.
Let $L=I_N\langle s\rangle$, $s\geq0$, $N\geq1$,
and $W=N+s\geq4$.
Assume $W$ is above both targets' horizontal and
vertical cutoffs through form degree four.
Let $m\ne0$ be an integer and suppose given actual
simplicial $S_R$-linear maps
\[
 z:L\otimes_R K(2)\to Z^4(X),\qquad
 \eta:L\otimes_R K(2)\to X^3,\qquad d\eta=mz.
 \tag{AT.74}
\]
Then there is an explicitly constructed strict PD
extension $X\hookrightarrow Y^{\mathrm{tf}}$ with
actual maps to $A,B$ extending the given maps,
such that every element in the image of $z$ becomes
a vertical boundary.
The terms of $Y^{\mathrm{tf}}$ are flat over $R$.
All weights below $W$ are unchanged, and the entire
simplex-zero and simplex-one algebras are unchanged.
The rationalized inclusion is a quasi-isomorphism
in every simplex and weight, and
$H^0((Y^{\mathrm{tf}})_a)=S_R$.

No vanishing of its positive integral cohomology,
and in particular no assertion about its nonlinear
cohomology in degrees one through three, is included.
\end{theorem}

\begin{proof}
Choose a graded free resolution $Q_\bullet\to L$
generated in weights at least $W$, by the explicit
weightwise construction in AT11.
Every term and every kernel has finite free
$R$-weight groups. Write $\varepsilon:Q_0\to L$
for the augmentation and $\partial$ for the
resolution differential, with $\partial Q_0=0$
inside its unaugmented complex.

Tensor this resolution with the two-term generator
complex $Rv\to Ru$, with $v$ in form degree two,
$u$ in form degree three, and differential $m$.
The following sign convention fixes the resulting
cochain complex $P$:
\[
 P^i=Q_{3-i}u\oplus Q_{2-i}v,\qquad
 d_P(a,b)=(-\partial a+mb,\partial b).
 \tag{AT.75}
\]
Terms with a negative $Q$ index are zero.
The two mixed terms cancel in $d_P^2$, so this
is a complex. It maps onto
$L[2]\xrightarrow{m}L[3]$ by the augmentation,
and that map is a quasi-isomorphism, since it
is the cone of multiplication by $m$ on a free
resolution of $L$.
Multiplication by $m$ is injective on $L$.
Thus $H^3(P)=L/mL$ and all other cohomology is zero.

Truncate without changing the last kernel:
\[
 \begin{aligned}
 W^3&=Q_0u,&
 W^2&=Q_1u\oplus Q_0v,&
 W^1&=Q_2u\oplus Q_1v,\\
 W^0&=\operatorname{coker}(P^{-1}\to P^0)
       =\ker(P^1\to P^2).
 \end{aligned}
 \tag{AT.76}
\]
The last equality is the canonical isomorphism
induced by $d_P:P^0\to P^1$; it follows from
$H^0(P)=H^1(P)=0$.
The degree-zero differential is the inclusion
of that kernel, and the other differentials
are induced by \textup{(AT.75)}.
Every weight of $W^0$ is free over $R$, since
it is a submodule of the finite free weight
group of $P^1$.
It need not be free over $S_R$.

There is a completely specified degree-one
twisting map
$\delta:W\otimes_R K(2)\to X$:
\[
 \delta(u_0)=z\varepsilon,\qquad
 \delta(v_0)=-\eta\varepsilon,\qquad
 \delta(u_i)=\delta(v_i)=0\ (i>0),\qquad
 \delta(W^0)=0.
 \tag{AT.77}
\]
The notation indicates its restrictions to the
corresponding $Q_i$ summands and the common
$K(2)$ parameter.
It obeys $d_X\delta+\delta d_W=0$.
For $v_0$ this is precisely $-d\eta+mz=0$;
for the other potentially nonzero terms it is
$\varepsilon\partial=0$.
The original twist on $P^0$ was zero, so it
indeed descends to the quotient in (AT.76).

Form the actual graded algebra
\[
 Y=X\otimes_{S_R}
       \operatorname{FreePD}_{S_R}(W\otimes_R K(2))
 \tag{AT.78}
\]
and define its differential on a new generator $t$
by $d_Yt=d_Wt+\delta(t)$.
The cocycle identity just checked proves its square
is zero. The derivation and PD differentiation
rules extend it through the free presentation.
Every relation is $S_R$-linear and homogeneous,
and all formulas are simplicial.
The graded augmentation of the free factor is
a left inverse to $X\to Y$, proving injectivity.
Since $Q_0\to L$ is surjective, the new $u_0$
generators supply primitives of every value of $z$.

We now construct the target maps.
Take $C_A$ from \textup{(AT.72)}.
If $F_r=P^{3-r}$, then $F_\bullet$ is a free
graded resolution of $L/mL$.
The twist \textup{(AT.77)} is a degree-four
cocycle in the total complex
$\operatorname{Hom}_{\mathrm{gr}S_R}(F_\bullet,C_A)$.
It has the $u_0$ component in resolution degree
zero, cochain degree four, and the $v_0$ component
in resolution degree one, cochain degree three.
The comparison \textup{(AT.73)} and the fact
that every free generator has weight at least $W$
give, on the finite diagonal of total degree four,
\[
 H^4\operatorname{Tot}\operatorname{Hom}(F_\bullet,C_A)
       \cong
 \operatorname{Ext}^2_{\mathrm{gr}S_R}(L/mL,S_R)=0.
 \tag{AT.79}
\]
For clarity, the vertical-cohomology spectral sequence
has only the row $j=2$ in this range, where its
coefficient Hom complex is
$\operatorname{Hom}(F_\bullet,S_R)$.
All terms entering or leaving the displayed total
degree use this same finite range.
The last vanishing is AT07 over $\mathbb Z$,
or the explicit local injective resolution
\textup{(AT.45)} over $\mathbb Z_{(p)}$.
The source is nonnegatively graded, exactly the
category in which those resolutions are injective.
No projectivity of the sphere parameter is asserted.

Choose a degree-three cochain whose differential
is that twist. In terms of the cochain indexing
of $P$, its components are actual simplicial maps
$h_i:P^i\otimes_R K(2)\to A^i$ for $0\leq i\leq3$,
satisfying
$d_Ah-hd_P=f\delta$.
Its terminal equation on $P^{-1}$ says
$h_0d_P=0$, so $h_0$ factors through $W^0$.
Thus it gives actual generator values on all
of $W\otimes K(2)$ satisfying precisely the
differential in \textup{(AT.78)}.
They define a strict PD map $Y\to A$.
The same calculation for $B$ gives its map.

Here the rational comparison has a particularly
explicit contraction, so no preservation theorem
for a nonlinear functor is being presumed.
On $P_{\mathbb Q}$ put
\[
                   H(u_i)=v_i/m,\qquad H(v_i)=0.
 \tag{AT.80}
\]
The sign in \textup{(AT.75)} gives
$d_PH+Hd_P=1$.
It descends to $W_{\mathbb Q}$: in degree one
take its value in $P^0$ modulo $\operatorname{im}P^{-1}$,
and in degree zero put it equal to zero.
The degree-zero contraction identity follows
from the same identity on $P^0$ modulo that image.
This is an actual weight-preserving $S_{\mathbb Q}$-linear
contraction, commuting with the $K(2)$ parameter.

Set $a=-\delta H$ on the rational generators.
The cocycle identity implies
$d_Xa-ad_W=\delta$.
The invertible change of generators $t'=t-a(t)$,
with inverse $t=t'+a(t)$, therefore untwists the
differential in \textup{(AT.78)}.
It leaves the free PD algebra on the contractible
complex $W_{\mathbb Q}\otimes K(2)$.
Extend $H$ to its graded derivation.
Its commutator with the differential is the
arity operator, because that identity holds on
every linear generator.
On positive arity $r$, dividing that derivation
by $r$ gives a contraction. These rational
operations are defined on each finite monomial
and preserve weight and the parameter's simplicial
maps. Tensoring the contraction with $X_{\mathbb Q}$
proves that the actual inclusion is a rational
quasi-isomorphism.

Finally put $Y^{\mathrm{tf}}=\operatorname{im}
(Y\to Y\otimes_R\mathbb Q)$.
The torsion ideal is differential, simplicial,
and PD-stable as in AT13.
Both maps to the torsionfree targets annihilate it.
The original $X$ still embeds, since its graded
retraction was torsionfree.
The quotient has flat terms over the PID $R$.
All new generators have weights at least $W$,
and $K(2)$ is zero in simplices zero and one.
The corresponding lower weights and simplex
algebras are therefore literally unchanged.
A closed degree-zero element in a changed weight
is, by the rational comparison, a rational
coefficient; its actual image in $A$ is an
integral coefficient by the chosen vertical cutoff.
As in AT13, injectivity on rational coefficients
forces the original element to belong to $S_R$.
Thus $H^0(Y^{\mathrm{tf}})=S_R$.
All the stated claims now follow.
\end{proof}

The calculation \textup{(AT.79)} settles the actual
linear target lifting for this construction.
It does not identify the low integral cohomology
of the nonlinear algebra in \textup{(AT.78)}.
The bottom module in \textup{(AT.76)} has genuine
$S_R$-relations. A splitting of its underlying
$R$-complex is not an $S_R$-linear PD comparison
with a free collection of disks.
No such low-cohomology comparison, and no resulting vanishing
of nonlinear $H^1,H^2,H^3$, is inferred here.

\paragraph{AT16-comparison (Proved; the actual map to the two-term tail presentation).}
Let $Z$ be the torsionfree quotient of the free
relative PD algebra on
$(L\otimes K(2))[2]\oplus(L\otimes K(2))[3]$,
with generators $v(l),u(l)$ and differential
$du(l)=z(l)$, $dv(l)=mu(l)-\eta(l)$.
This has the two-term presentation of A-S09-21,
now for the specified maps $z,\eta$ and sphere parameter.
The integral cohomology assertion of A-S09-21 retains
the specific independent-class hypotheses stated there.
There is an actual surjective strict PD-DGA map
over $X$
\[
                 Y^{\mathrm{tf}}\longrightarrow Z
 \tag{AT.81}
\]
defined on the resolved generators by
\[
 u_0(a)\longmapsto u(\varepsilon a),\qquad
 v_0(a)\longmapsto v(\varepsilon a),\qquad
 u_i,v_i\longmapsto0\ (i>0),\qquad W^0\longmapsto0.
 \tag{AT.82}
\]
Indeed the $u_0,v_0$ differential equations are
exactly those defining $Z$. The other equations
map to zero by $\varepsilon\partial=0$;
the bottom differential lands in $W^1$, which
is also sent to zero.
The free PD universal property gives the map
before the torsion quotient, and the torsionfree
target makes it factor through $Y^{\mathrm{tf}}$.
Its images generate $Z$ because $\varepsilon$
is surjective, proving surjectivity.
The construction is simplicial and preserves
both gradings and the given map from $X$.

This map is a rational quasi-isomorphism.
For $Z_{\mathbb Q}$ the change
$u(l)\mapsto u(l)-\eta(l)/m$ gives the free
contractible two-term disk on $L_{\mathbb Q}\otimes K(2)$.
Both its coefficient inclusion and that into
$Y_{\mathbb Q}$ are rational quasi-isomorphisms,
and their composite with (AT.81) is the same
actual inclusion of $X_{\mathbb Q}$.
It follows on cohomology that (AT.81) is an
isomorphism after rationalization.
Already on the linear generator complexes,
$W\to(L[2]\xrightarrow{m}L[3])$ is an integral
quasi-isomorphism, by the calculation of $H(P)$
and the exact truncation in (AT.76).
No integral quasi-isomorphism of their nonlinear
free PD algebras, or factorization of the maps
to $A,B$ through (AT.81), is inferred from this.


